\documentclass[compress, 8pt]{beamer}

\usepackage{presentationtemplate}
\usepackage[askip=3mm, bskip=3mm]{terminal}
\usepackage[linenosfontsize=\tiny, askip=3mm, bskip=3mm]{mylisting}
\tcbuselibrary{hooks}

\newtcolorbox{task}{
    colback=yellow!50!white,
    boxrule=0.02cm,
    colframe=black,
    sharp corners,
    left=0mm,
    right=0mm,
    top=0mm,
    bottom=0mm,
    before upper={\textbf{Задание}:\:},
}

\title{Наследование}

\begin{document}

  \frame[plain]{\titlepage}

  \begin{frame}[fragile]

    \frametitle{Наследование: синтаксис}

    \begin{itemize}
      \item Класс, образованный из другого класса, называется
        \textbf{производным} (derived); исходный \textemdash \space
        \textbf{базовым} (base).
      \item Синтаксис: \texttt{struct Derived : Base \{\};} \textemdash \space
        поля базового класса доступны производному через точку.
    \end{itemize}

    \myinputlisting[minted language=cpp]
        {Presentations/12-Inheritance/}
        {syntax.cpp}

  \end{frame}

  \begin{frame}[fragile]

    \frametitle{Спецификаторы наследования}

    По умолчанию наследование \texttt{public} у \texttt{struct}
    и \texttt{private} у \texttt{class}. Спецификатор можно указать явно.

    \myinputlisting[minted language=cpp]
        {Presentations/12-Inheritance/}
        {access-specs.cpp}

    \vspace{1em}

    \begin{center}
      \begin{tabular}{lll}
        \hline
        Член \texttt{Base}
          & \texttt{D1}, \texttt{D2}, \texttt{D4} (\texttt{public})
          & \texttt{D3} (\texttt{private}) \\
        \hline
        \texttt{public} & \texttt{public} & \texttt{private} \\
        \texttt{protected} & \texttt{protected} & \texttt{private} \\
        \texttt{private} & недоступен & недоступен \\
        \hline
      \end{tabular}
    \end{center}

  \end{frame}

  \begin{frame}[fragile]

    \frametitle{\texttt{protected}-члены}

    \texttt{protected}-члены доступны методам производного класса,
    но недоступны из любого другого кода.

    \myinputlisting[minted language=cpp, minted options app={fontsize=\scriptsize}]
        {Presentations/12-Inheritance/}
        {protected.cpp}

    \begin{terminalwindow}
!\shellcommand{./a.out}!
a = 1, b = 2
    \end{terminalwindow}

  \end{frame}

  \begin{frame}[fragile]

    \frametitle{Задача: найдите ошибку}

    \begin{task}
      Найдите ошибку в коде ниже.
    \end{task}

    \myinputlisting[minted language=cpp]
        {Presentations/12-Inheritance/}
        {task2.cpp}

  \end{frame}

  \begin{frame}[fragile]

    \frametitle{Конструкторы базового класса}

    \begin{itemize}
      \item Конструкторы базового класса НЕ наследуются: вызова
        \texttt{Derived(1)}, который напрямую вызвал бы
        \texttt{Base(int)}, не существует.
      \item Конструктор базового класса вызывается явно
        \textemdash \space в инициализаторном списке конструктора
        производного класса.
    \end{itemize}

    \myinputlisting[minted language=cpp]
        {Presentations/12-Inheritance/}
        {base-ctors.cpp}

  \end{frame}

  \begin{frame}[fragile]

    \frametitle{Задача: какой конструктор базового класса будет вызван}

    \begin{task}
      Подумайте, что выведет программа. Почему?
    \end{task}

    \myinputlisting[minted language=cpp]
        {Presentations/12-Inheritance/}
        {task3.cpp}

  \end{frame}

  \begin{frame}[fragile]

    \frametitle{Порядок вызова конструкторов и деструкторов}

    Конструкторы: сначала базовый, потом производный. Деструкторы
    \textemdash \space в обратном порядке: сначала производный,
    потом базовый.

    \myinputlisting[minted language=cpp]
        {Presentations/12-Inheritance/}
        {ctor-order.cpp}

  \end{frame}

  \begin{frame}[fragile]

    \frametitle{Порядок вызова конструкторов и деструкторов}

    Вывод программы с прошлого слайда:

    \begin{terminalwindow}
!\shellcommand{./a.out}!
Base
Derived
~Derived
~Base
    \end{terminalwindow}

  \end{frame}

  \begin{frame}[fragile]

    \frametitle{Задача: порядок вызовов}

    \begin{task}
      Подумайте, что выведет программа. Обоснуйте ответ.
    \end{task}

    \myinputlisting[minted language=cpp]
        {Presentations/12-Inheritance/}
        {task1.cpp}

  \end{frame}

  \begin{frame}[fragile]

    \frametitle{Object slicing}

    При копировании, присваивании или передаче объекта производного
    класса по значению копируется только под-объект базовой части.
    Это называется \textbf{обрезание объекта} (object slicing).

    \myinputlisting[minted language=cpp]
        {Presentations/12-Inheritance/}
        {slicing.cpp}

  \end{frame}

  \begin{frame}[fragile]

    \frametitle{Указатель на базовый класс и расположение объекта}

    Базовый под-объект размещён в начале объекта производного класса:
    адрес объекта и адрес его базовой части совпадают. Поэтому
    указатель на базовый класс может указывать на производный объект.

    \begin{columns}[T]

      \begin{column}{0.48\textwidth}

        \myinputlisting[minted language=cpp]
            {Presentations/12-Inheritance/}
            {layout.cpp}

      \end{column}

      \begin{column}{0.48\textwidth}

        \begin{terminalwindow}
!\shellcommand{g++ layout.cpp}!
!\shellcommand{./a.out}!
0x7ffc0ef376e8
0x7ffc0ef376e8
0x7ffc0ef376ec
        \end{terminalwindow}

      \end{column}

    \end{columns}

  \end{frame}

  \begin{frame}[fragile]

    \frametitle{Скрытие метода}

    \begin{itemize}
      \item Метод производного класса с тем же именем не переопределяет,
        а \textbf{скрывает} метод базового класса.
      \item Без \texttt{virtual} вызывается метод, соответствующий
        \textbf{статическому} типу выражения. Как сделать, чтобы
        вызывался метод динамического типа, \textemdash \space
        в следующей лекции.
    \end{itemize}

    \myinputlisting[minted language=cpp]
        {Presentations/12-Inheritance/}
        {hiding.cpp}

  \end{frame}

  \begin{frame}[fragile]

    \frametitle{\texttt{using Base::}}

    \begin{itemize}
      \item Заявление \texttt{using Base::Foo;} вводит в область видимости
        производного класса все перегрузки базового класса.
      \item Без него объявленный в производном классе метод скрывает
        ВСЕ базовые перегрузки: \texttt{d.Foo(1)} молча вызовет
        \texttt{Foo(double)}, а \texttt{d.Foo()} \textemdash \space
        ошибка компиляции.
    \end{itemize}

    \myinputlisting[minted language=cpp, minted options app={fontsize=\scriptsize}]
        {Presentations/12-Inheritance/}
        {using-base.cpp}

  \end{frame}

\end{document}
